\documentclass[11pt]{article}

\usepackage{amsmath,amssymb}

\title{Pattern Recognition and Machine Learning Study Notes}
\author{Guilherme Candinho}
\date{\today}

\begin{document}

\maketitle

\section{Introduction}
\subsection{Definitions}
\begin{itemize}
  \item Supervised learning problems are applications in which the training data comprises examples of the input vectors along with their corresponding target vectors.
  \item Classification tasks are cases where the aim is to assign each input vector to one of a finite number of discrete categories.
  \item Regression tasks are cases where the aim is to assign each input vector to an output that consists of one or more continuous variables.
  \item Unsupervised learning problems are application in which the training data does not have any corresponding target values. The goal is to may discover groups of similar examples within the data, where it is called clustering, or to determine the distribution of data within the input space, know as density estimation, or to project the data from a high-dimensional space down to two or three dimension for the purpose of visualization.
  \item Reinforcement learning is concerned with problems of finding suitable actions to take in a given situation in order to maximize a reward. This is achieved by a process of trial and error.
\end{itemize}
Although each of these tasks need its own tools and techniques, many of the key ideas that underpin them are commom to all such problems. The first chapter, gives an introduction to several of these tools and techniques therefore, this notes will have some of the examples and my toughs on such problems.
\subsection{Polynomial Curve}
In general, when working with real world data, whe may wish to find an underlying regularity on the data to be able to assign new observations correctly. Such real world data is likely to allways have some kind of noises due to intrinsically stochastic (i.e random) processes or due to being some source of variability that are unobserved.


\end{document}
